\section{Related work}
\subsection{Aerial embodied agents and simulation platforms}
Vision-and-language navigation connects natural-language instructions to sequential visual decisions~\cite{anderson2018vision}. AerialVLN extends this problem to UAV navigation, while CityNav introduces large-scale real-world aerial navigation data~\cite{liu2023aerialvln,lee2025citynav}. OpenFly provides a broader aerial vision-language platform, toolchain, and trajectories across multiple heights and scenes~\cite{gao2026openfly}. These works establish language-conditioned aerial motion and height-varying observation as active research areas. DisasterClaw addresses a complementary requirement: constructing task-scoped, paired disaster observations and measuring whether an inspection action changes the requested disaster assessment.

Large pretrained models can be composed with perception and navigation modules without treating the entire stack as a newly learned end-to-end policy~\cite{shah2023lmnav}. CityNavAgent uses hierarchical semantic planning and global memory, spatial text representations support language-model reasoning in aerial navigation, and AerialClaw exposes modular aerial skills and feedback~\cite{zhang2025citynavagent,gao2024stmr,li2026aerialclaw}. DisasterClaw adopts this modular view. Semantic maps, planning, memory, and tool execution support the implemented agent; the platform contribution lies in their integration with georeferenced disaster scenes, controllable inspection observations, and a common measurement interface.

CityEQA connects active urban exploration to embodied question answering and evaluates answers after view refinement~\cite{zhao2025cityeqa}. Its setting is a direct precedent for active aerial question answering. DisasterClaw differs in task domain, source imagery, action abstraction, and evaluation target. Rather than benchmarking general-purpose flight skills or high-fidelity navigation, it couples a georeferenced disaster scene to a controlled observation transition and measures whether that transition changes task-relevant evidence and the terminal inspection answer. This narrower scope motivates reporting observation construction, action execution, and task-level correction or harm with equal precision.

Recent work has expanded aerial embodied intelligence beyond instruction-following navigation. AeroVerse organizes UAV-agent evaluation around scene awareness, spatial reasoning, navigational exploration, task planning, and motion decision, providing a broad benchmark suite for aerospace embodied world models~\cite{yao2024aeroverse}. BEDI further formalizes UAV embodied-agent evaluation through a perception–decision–action chain and evaluates semantic and spatial perception, motion control, tool use, planning, and action generation across real and virtual settings~\cite{guo2026bedi}. OpenUAV and the UAV-Need-Help benchmark emphasize realistic aerial trajectory execution and assistant-guided object search in three-dimensional environments~\cite{wang2024openuav}. These platforms address general aerial embodiment, realistic navigation, and standardized skill evaluation. DisasterClaw is complementary rather than competing in simulator fidelity: it uses georeferenced pre-/post-disaster imagery to make the task value of a controlled view change directly observable at the final disaster-inspection answer.

\subsection{Disaster interpretation from remote-sensing imagery}
xBD provides paired satellite images and building-damage annotations for disaster assessment~\cite{gupta2019xbd}. RescueADI studies autonomous tool use for adaptive interpretation of remote-sensing disaster images~\cite{liu2024rescueadi}. Thus, paired damage recognition and agentic disaster analysis both predate this work. Recent remote-sensing resources further broaden the evaluation of multimodal reasoning and agentic analysis. ThinkGeo evaluates tool-augmented language-model agents on multi-step remote-sensing tasks, including disaster assessment and change analysis, and scores both intermediate tool execution and final answers~\cite{thinkgeo2025}. DisasterM3 provides a large-scale multi-hazard, multi-sensor vision-language benchmark for disaster damage assessment and response, covering perception and reasoning across diverse disaster events~\cite{wang2025disasterm3}. These resources strengthen static-image reasoning, multimodal generalization, and structured tool use. DisasterClaw addresses a different axis: the observation itself is an environment variable, so the agent can execute a view-changing action and the benchmark can measure whether the newly acquired evidence improves or harms the final task answer. DisasterClaw uses such static records as the substrate of an interactive inspection environment: a geographic window becomes an observation, vehicle state changes the window, and the resulting episode is scored at both building and task levels.

The distinction is a data-to-world construction route: registered imagery defines an observation environment, while source annotations and geographic rules define the task benchmark. This is complementary to the navigation, tool-use, and general embodiment objectives above, not a claim that all other platforms use only synthetic scenes. The transformation remains constrained by the data. Cropping and resampling an orthographic product can alter the number of pixels assigned to a target, but it cannot synthesize unrecorded viewpoints or detail below the native source scale. The platform is consequently suited to controlled studies of spatial coverage, observation allocation, and answer formation. Evidence from it cannot be presented as a measurement of real UAV image formation or field performance.

\subsection{Active perception, uncertainty, and task value}
Active perception treats sensing as an action selected in service of a task~\cite{bajcsy1988active}. Planning under partial observability further connects information gathering with uncertain states and action outcomes~\cite{curtis2024tampura}. This perspective motivates the central platform requirement: a new observation should be evaluated by the decision it enables and the cost it incurs.

Entropy reduction is a convenient trigger but is not equivalent to expected task utility. A distribution can become more concentrated around an incorrect class, while a view may help a count or spatial relation without changing the most uncertain building label. DisasterClaw therefore records classification, correction and harm, final question correctness, and executed observation count as separate endpoints.

Temperature scaling addresses probability calibration, while evaluation under dataset shift shows why in-distribution calibration should not be treated as a universal reliability property~\cite{guo2017calibration,ovadia2019can}. Selective prediction measures the trade-off between coverage and error, and uncertainty-aware language-model planning has been studied through help-seeking behavior~\cite{geifman2019selectivenet,ren2023knowno}. We use these distinctions to separate raw and calibrated binary-damage triggers, prediction-set sensing rules, and the VLM's generated answer. None alone establishes reliability for the composed agent.
